% STATUS: draft
% LAST_MODIFIED: 2026-07-26
\documentclass{internal-report}
\usepackage[UTF8]{ctex}

\title{2020 CUMCM C 题 — 初始分析报告}
\author{MathModel Team}
\date{2026年7月26日}

\begin{document}
\maketitle

\section{问题重述}

某银行需对中小微企业进行信贷决策：评估信贷风险、决定是否放贷、确定贷款额度和利率。
数据包括 123 家有信贷记录企业（附件1，含信誉评级和违约标签）和 302 家无信贷记录企业
（附件2，仅有发票数据），以及利率与客户流失率的关系数据（附件3）。

三个子问题：
\begin{enumerate}
  \item 对 123 家企业量化信贷风险，在年度信贷总额固定时给出信贷策略。
  \item 对 302 家企业量化信贷风险，在年度总额 1 亿元时给出信贷策略。
  \item 考虑突发事件（如新冠疫情）的差异化影响，给出信贷调整策略。
\end{enumerate}

\section{子问题拆解}

\subsection{子问题 1：有标签企业的风险评估（附件1）}

\textbf{输入}：123 家企业的进项/销项发票数据、信誉评级（A/B/C/D）、违约标签。

\textbf{输出}：每家企业违约概率的量化估计 + 信贷决策（额度、利率）。

\textbf{核心挑战}：从发票数据中提取有效风险特征。文献中使用了以下几种方案：
\begin{itemize}
  \item 11 指标 + 熵权 + TOPSIS（D1）—— 客观赋权，但忽略了指标相关性
  \item 6 指标 + Logistic 回归（D2）—— 可解释但准确率有限（84.6\%，违约召回率仅 48.1\%）
  \item 20 指标 + Soft Voting 集成（D3）—— AUC=0.994，但计算复杂
  \item 7 指标 + PCA 降维 → 安全指数（D4）—— 含 FFT 交易规律分析，创新性强
\end{itemize}

\subsection{子问题 2：无标签企业的风险评估（附件2）}

\textbf{核心挑战}：附件 2 的 302 家企业没有信誉评级和违约标签。
需要用附件 1 训练的模型推断附件 2 的风险。

文献中的方案：
\begin{itemize}
  \item D2：用欧氏距离匹配——将附件 2 企业的特征向量与附件 1 企业比较，
    以最近邻的信用等级作为预测等级，再代入 Logistic 模型。
  \item D3：Soft Voting + XGBoost 直接预测，准确率 87.8\%。
  \item D4：BP 神经网络预测风险。
\end{itemize}

\textbf{年度总额 1 亿的分配}：所有文献都采用了多目标规划——在总贷款额度约束下，
最大化风险调整后收益，最小化最大单笔风险和客户流失。

\subsection{子问题 3：突发因素调整}

\textbf{核心挑战}：量化新冠疫情对每个企业的影响，但没有精确的「冲击值」数据。

文献方案：
\begin{itemize}
  \item D3：系统聚类 → 将 302 家企业按疫情表现分为 3 类 → 每类赋风险乘数
    $Q_1=1.2, Q_2=1.5, Q_3=2.0$ → 修正违约概率为 $(1+Q_k)\theta$。
  \item D1：三级因子法（固定因子 + 变量因子 + 模型因子），需要行业统计数据。
  \item D4：经验评估 + 比例因子。
\end{itemize}

\section{相关知识}

\subsection{风险量化方法选型}

从知识库中识别出四种建模哲学（SP-005）：

\begin{table}[h]
\centering
\caption{四种方案对比}
\begin{tabular}{lcccc}
\toprule
维度 & A: TOPSIS & B: Logistic & C: 集成学习 & D: PCA+GA \\
\midrule
指标数 & 11 & 6 & 20 & 7 \\
核心理念 & 客观排序 & 概率模型 & 多模型集成 & 降维+优化 \\
创新点 & RAROC理论 & 欧氏距离匹配 & Voting AUC=0.994 & FFT交易规律 \\
决策模型 & RAROC最大化 & 多目标 & 多目标 & GA非线性 \\
突发事件 & 三级因子 & 指标修正 & 聚类+乘数 & 比例因子 \\
\bottomrule
\end{tabular}
\end{table}

\subsection{关键陷阱（PF-007）}

类别不平衡问题：123 家中仅约 27 家违约（22\%）。
仅看总体准确率会严重误导——D1 的 Logistic 回归准确率 84.6\%，
但违约类召回率仅 48.1\%，漏掉了超过一半的真正违约者。

\textbf{必须同时报告混淆矩阵、召回率和 AUC}，不能仅依赖总体准确率。

\subsection{信贷决策模型（MS-025）}

三篇论文的共同框架：多目标规划。
目标 1：最大化风险调整后的总收益。
目标 2：最小化最大单笔风险（minimax 准则）。
目标 3：最小化客户流失率（由利率驱动，见附件 3）。

约束：总贷款额 ≤ 预算、单笔 10--100 万、年利率 4\%--15\%、
信誉评级 D 的企业原则上不放贷。

\section{初始建模假设}

\begin{enumerate}
  \item 企业间信贷风险相互独立（不考虑行业联动违约）。
  \item 当期决策相互独立（非动态规划，不考虑跨期依赖）。
  \item 利率与客户流失率的关系稳定（基于附件 3 的 2019 年数据）。
  \item 发票数据真实反映企业经营状况（有效发票→正常交易，作废/负数→异常）。
\end{enumerate}

\section{待解决的开放问题}

\begin{enumerate}
  \item \textbf{特征工程方案选择}：用 D1 的 11 指标、D2 的 6 指标，还是设计更优的特征集？
  \item \textbf{风险量化方法选型}：四种方案各有优劣，需要根据实际情况选择
    （可解释性 vs. 预测精度 vs. 创新性）。
  \item \textbf{问题 1 的年度信贷总额}：题目未明确，需要合理假设或从数据中反推。
  \item \textbf{附件 2 企业的行业识别}：企业名称被脱敏，如何从发票数据推断行业归属？
  \item \textbf{突发因素的冲击量化}：无外部宏观数据时，如何更客观地确定风险乘数？
  \item \textbf{模型验证}：附件 1 仅有 123 个样本，训练-测试切分导致方差大。
    是否需要留出法、交叉验证或 Bootstrap？
  \item \textbf{利率-流失率的函数形式}：附件 3 的数据结构未确认。
    需要选择线性/对数/多项式拟合。
\end{enumerate}

\end{document}
